\documentclass[11pt]{article}

\usepackage[letterpaper,margin=0.72in,headheight=22pt,footskip=28pt]{geometry}
\usepackage[T1]{fontenc}
\usepackage{lmodern}
\usepackage{microtype}
\usepackage[table]{xcolor}
\usepackage{amsmath,amssymb}
\usepackage{array,booktabs,tabularx}
\usepackage{enumitem}
\usepackage{fancyhdr}
\usepackage{lastpage}
\usepackage[hidelinks]{hyperref}

\definecolor{navy}{HTML}{173B4D}
\definecolor{bluegray}{HTML}{EAF4F7}
\definecolor{warmgray}{HTML}{F5F3EE}
\definecolor{rulegray}{HTML}{C9CED1}

\setlength{\parindent}{0pt}
\setlength{\parskip}{5pt}
\setlist[enumerate]{leftmargin=1.55em,itemsep=4pt,topsep=3pt}
\setlist[itemize]{leftmargin=1.45em,itemsep=2pt,topsep=3pt}
\renewcommand{\arraystretch}{1.18}

\pagestyle{fancy}
\fancyhf{}
\lhead{\small CSCI 340}
\rhead{\small Assignment 2}
\cfoot{\small Page \thepage\ of \pageref{LastPage}}
\renewcommand{\headrulewidth}{0.4pt}
\renewcommand{\headrule}{\hbox to\headwidth{\color{rulegray}\leaders\hrule height \headrulewidth\hfill}}

\newcommand{\sectionbar}[2]{%
  \vspace{4pt}%
  {\color{navy}\rule{\linewidth}{1.3pt}}\vspace{2pt}
  {\Large\bfseries #1}\hfill{\large\bfseries #2}\par
  {\color{navy}\rule{\linewidth}{0.5pt}}\vspace{3pt}
}
\newcommand{\question}[3]{%
  \vspace{5pt}\noindent{\large\bfseries #1\quad #2}\hfill{\bfseries #3 points}\par\vspace{1pt}
}
\newcommand{\anspace}[1]{\par\vspace{#1}}
\newcommand{\schema}[1]{\texttt{#1}}
\newcommand{\student}{\mathsf{student}}
\newcommand{\course}{\mathsf{course}}
\newcommand{\enrollment}{\mathsf{enrollment}}

\begin{document}

{\small\bfseries CSCI 340 \quad DATABASE DESIGN}\par
\vspace{4pt}
{\fontsize{22}{25}\selectfont\bfseries Relational Algebra Practice}\par
{\large Selection, projection, joins, and set operations}\par
\vspace{9pt}

\begin{tabularx}{\linewidth}{@{}>{\bfseries}l X >{\bfseries}l X@{}}
Name & \hrulefill & People consulted & \hrulefill \\
\end{tabularx}

\vspace{8pt}
\textbf{Purpose.} Build fluency with relational algebra by reading expressions, tracing exact results, and writing expressions of increasing complexity. Every question uses the same small registration database.

\textbf{Directions.} Submit your own completed worksheet. You may work with classmates, but list everyone you worked with above. Unless a question says otherwise, give only the requested attributes. When asked to trace an expression, write the output schema and every resulting tuple. Tuple order does not matter because a relation is an unordered set.


\subsection*{Notation and assumptions}

\begin{tabularx}{\linewidth}{@{}>{\bfseries}p{1.15in} >{\centering\arraybackslash}p{1.25in} X@{}}
\toprule
Operation & Form & Effect \\
\midrule
Selection & $\sigma_p(E)$ & Keeps tuples of $E$ that satisfy predicate $p$. \\
Projection & $\pi_{A_1,\ldots,A_k}(E)$ & Keeps the listed attributes and removes duplicate output tuples. \\
Natural join & $E_1 \bowtie E_2$ & Matches on every shared attribute name and keeps one copy of each shared attribute. \\
Theta join & $E_1 \bowtie_{\theta} E_2$ & Matches pairs satisfying $\theta$ and keeps attributes from both inputs. \\
Set operations & $E_1\cup E_2$, $E_1\cap E_2$, $E_1-E_2$ & Combine union-compatible inputs as sets of tuples. \\
\bottomrule
\end{tabularx}

Use qualified names such as $\student.\mathit{student\_id}$ when an expression could otherwise be ambiguous. You may use assignment notation, $x\leftarrow E$, to name an intermediate result. String values appear in quotation marks.

\newpage
\sectionbar{Shared registration database}{Reference for every section}

The schema has three relations:

\begin{align*}
\student(&\mathit{student\_id},\ \mathit{student\_name},\ \mathit{major},\ \mathit{year})\\
\course(&\mathit{course\_id},\ \mathit{course\_title},\ \mathit{department},\ \mathit{credits})\\
\enrollment(&\mathit{student\_id},\ \mathit{course\_id},\ \mathit{grade})
\end{align*}

\begin{minipage}[t]{0.49\linewidth}
\textbf{Relation: student}

\rowcolors{2}{bluegray}{white}
\small\setlength{\tabcolsep}{3.5pt}
\begin{tabular}{@{}rlll@{}}
\rowcolor{navy}\color{white}\textbf{student\_id} & \color{white}\textbf{student\_name} & \color{white}\textbf{major} & \color{white}\textbf{year} \\
101 & Ava & CS & junior \\
102 & Jordan & Math & sophomore \\
103 & Lin & Biology & senior \\
104 & Sam & CS & sophomore \\
105 & Noor & Math & junior \\
106 & Mira & History & first-year \\
\end{tabular}
\end{minipage}\hfill
\begin{minipage}[t]{0.49\linewidth}
\textbf{Relation: course}

\rowcolors{2}{bluegray}{white}
\small\setlength{\tabcolsep}{3.5pt}
\begin{tabular}{@{}rllr@{}}
\rowcolor{navy}\color{white}\textbf{course\_id} & \color{white}\textbf{course\_title} & \color{white}\textbf{department} & \color{white}\textbf{credits} \\
C10 & Databases & CS & 3 \\
C20 & Discrete Math & Math & 3 \\
C30 & Ecology & Biology & 4 \\
C40 & Data Ethics & CS & 3 \\
C50 & World History & History & 4 \\
\end{tabular}
\end{minipage}

\vspace{12pt}
\begin{minipage}[t]{0.49\linewidth}
\textbf{Relation: enrollment}

\rowcolors{2}{bluegray}{white}
\begin{tabular}{@{}rrl@{}}
\rowcolor{navy}\color{white}\textbf{student\_id} & \color{white}\textbf{course\_id} & \color{white}\textbf{grade} \\
101 & C10 & A \\
101 & C40 & B \\
102 & C20 & A \\
103 & C30 & B \\
104 & C10 & B \\
105 & C10 & A \\
\end{tabular}
\end{minipage}\hfill
\begin{minipage}[t]{0.47\linewidth}
\textbf{Facts worth noticing}

\begin{itemize}
  \item A student may have more than one enrollment.
  \item A course may have more than one enrolled student.
  \item A student's major need not match a course's department.
  \item A student or course may have no matching enrollment.
\end{itemize}
\end{minipage}


\newpage
\sectionbar{Section A}{Selection and projection \quad 12 points}

\question{A1}{Read expressions in plain English}{3}
For each expression, state in plain English exactly what its result contains.

\begin{enumerate}[label=\alph*.]
  \item $\sigma_{\mathit{major}=\text{``CS''}}(\student)$
  \item $\pi_{\mathit{student\_name},\mathit{year}}\bigl(\sigma_{\mathit{major}=\text{``Math''}}(\student)\bigr)$
  \item $\pi_{\mathit{department}}(\course)$
\end{enumerate}
\anspace{1.05in}

\question{A2}{Trace exact results}{5}
Write the output schema and every tuple produced by each expression. Explicitly identify any duplicate tuple removed by projection.

\begin{enumerate}[label=\alph*.]
  \item $\pi_{\mathit{major}}(\student)$
  \anspace{0.95in}
  \item $\pi_{\mathit{student\_id},\mathit{student\_name}}\bigl(\sigma_{\mathit{year}=\text{``sophomore''}}(\student)\bigr)$
  \anspace{0.95in}
\end{enumerate}

\question{A3}{Write expressions}{4}
Write one relational-algebra expression for each request.

\begin{enumerate}[label=\alph*.]
  \item Return the title of every 4-credit course. \hfill (2)
  \anspace{0.55in}
  \item Return the name and major of every junior student. \hfill (2)
  \anspace{0.55in}
\end{enumerate}

\newpage
\sectionbar{Section B}{Natural join and theta join \quad 17 points}

\question{B1}{Distinguish the joins}{3}
Answer each question about the expression $\student\bowtie\enrollment$.

\begin{enumerate}[label=\alph*.]
  \item Which shared attribute is used automatically as the matching condition?
  \item Why does that shared attribute appear only once in the result schema?
  \item If the two relations had no shared attribute names, what would their natural join produce?
\end{enumerate}
\anspace{0.65in}

\question{B2}{Trace a natural join}{6}
Write the output schema and every tuple in $\student\bowtie\enrollment$. Then name the student who does not appear and explain why. Tuple order does not matter.
\anspace{2.85in}

\question{B3}{Analyze a theta join}{4}
Consider the expression
\[
\student\ \bowtie_{\student.\mathit{major}=\course.\mathit{department}}\ \course.
\]

\begin{enumerate}[label=\alph*.]
  \item State what the expression means in plain English. \hfill (1)
  \item State its output schema. \hfill (1)
  \item How many output tuples are there? Briefly justify your count; do not list every tuple. \hfill (2)
\end{enumerate}
\anspace{1.25in}

\newpage
\question{B4}{Write composed join queries}{4}
Write one relational-algebra expression for each request. Natural join is appropriate because the linked relations share deliberately matching attribute names.

\begin{enumerate}[label=\alph*.]
  \item Return the names of students enrolled in course C10. \hfill (2)
  \anspace{0.9in}
  \item Return the titles of courses taken by Math majors. \hfill (2)
  \anspace{0.9in}
\end{enumerate}

\newpage
\sectionbar{Section C}{Set operations and composition \quad 16 points}

For C1--C3, use these named intermediate results:
\begin{align*}
\mathit{cs\_student} &\leftarrow
\pi_{\mathit{student\_id}}\bigl(\sigma_{\mathit{major}=\text{``CS''}}(\student)\bigr)\\
\mathit{sophomore} &\leftarrow
\pi_{\mathit{student\_id}}\bigl(\sigma_{\mathit{year}=\text{``sophomore''}}(\student)\bigr)
\end{align*}

\question{C1}{Check compatibility}{3}
Are $\mathit{cs\_student}$ and $\mathit{sophomore}$ union-compatible? State the schema of each result and explain briefly.
\anspace{0.75in}

\question{C2}{Read set expressions in plain English}{3}
State what each expression means in plain English. Be precise about the direction of difference.

\begin{enumerate}[label=\alph*.]
  \item $\mathit{cs\_student}\cup\mathit{sophomore}$
  \item $\mathit{cs\_student}\cap\mathit{sophomore}$
  \item $\mathit{cs\_student}-\mathit{sophomore}$
\end{enumerate}
\anspace{0.85in}

\question{C3}{Trace set results}{4}
First write every tuple in $\mathit{cs\_student}$ and $\mathit{sophomore}$. Then write every tuple in the union, intersection, and difference from C2. Explain why a student ID appearing in both inputs occurs only once in the union.
\anspace{1.8in}

\newpage
\sectionbar{Section C continued}{Composition with joins}

\question{C4}{Compose set operations with joins}{6}
First define the IDs of students enrolled in a CS-department course:
\[
\mathit{takes\_cs}\leftarrow
\pi_{\mathit{student\_id}}
\left(
\enrollment\bowtie\sigma_{\mathit{department}=\text{``CS''}}(\course)
\right).
\]
Using $\mathit{cs\_student}$ and $\mathit{takes\_cs}$, write an expression for each request. You do not need to enumerate the results.

\begin{enumerate}[label=\alph*.]
  \item IDs of students who major in CS or take a CS course. \hfill 
  \anspace{0.5in}
  \item IDs of students who major in CS and take a CS course. \hfill
  \anspace{0.5in}
  \item IDs of students who take a CS course but do not major in CS. \hfill 
  \anspace{0.5in}
\end{enumerate}

\end{document}
